% Corrected French for printed p.191; departures are append-only below.
\clearpage
\SourcePage{191}

\LayerTarget{tohoku:section:4.3}
\textbf{4.3. Cas d'un faisceau d'anneaux constant.}
\LayerTarget{tohoku:setup:constant-ring-sheaf}
Supposons que $\mathbf O$ soit le faisceau d'anneaux constant défini par un
anneau avec unité $O$. Soit $M$ un $O$-module, et soit $\mathbf M$ le
$\mathbf O$-Module constant qu'il définit. Pour tout $\mathbf O$-Module $A$,
on a alors un isomorphisme naturel
\LayerTarget{tohoku:eq:4.3.1}
\[
\tag{4.3.1}
\mathrm{Hom}_{\mathbf O}(\mathbf M,A)=\mathrm{Hom}_O(M,\Gamma(A)).
\]
C'est là un isomorphisme fonctoriel, qu'on peut aussi écrire
\LayerTarget{tohoku:eq:4.3.2}
\[
\tag{4.3.2}
h_{\mathbf M}=h_M\Gamma
\]
\LayerTarget{tohoku:setup:functors-h-constant-module}
où on pose, comme au numéro précédent,
$h_{\mathbf M}(A)=\mathrm{Hom}_{\mathbf O}(\mathbf M,A)$,
$h_M(N)=\mathrm{Hom}_O(M,N)$ pour tout $\mathbf O$-Module $A$ et tout
$O$-module $N$. $\Gamma$ est considéré comme un foncteur de
$\mathbf C^{\mathbf O}$ dans la catégorie $\mathbf C^O$ des $O$-modules à gauche, et
$h_M$ comme un foncteur de $\mathbf C^O$ dans la catégorie des
groupes abéliens. Le théorème 2.4.1 s'applique, grâce au

\LayerTarget{tohoku:lemma:4.3:gamma-injective}
\textsc{Lemme} \textit{Si $A$ est un $\mathbf O$-Module injectif,
alors $\Gamma(A)$ est un $O$-module injectif.}

\LayerTarget{tohoku:proof:lemma-4.3-gamma-injective}
Il faut montrer que le foncteur
$M\to\mathrm{Hom}_O(M,\Gamma(A))$ est exact, ce qui résulte aussitôt
de la formule (4.3.1), puisque $M\to\mathbf M$ est un foncteur exact.
On obtient donc, grâce au théorème 2.4.1, et en remarquant que les foncteurs
dérivés de $\Gamma$ sont les mêmes, qu'on considère $\Gamma$ comme prenant
ses valeurs dans $\mathbf C^O$ ou dans la catégorie $\mathbf C$ des groupes abéliens :

\LayerTarget{tohoku:theorem:4.3.1}
\textsc{Théorème 4.3.1.} \textit{Soient $X$ un espace, $O$ un anneau
avec unité, $M$ un $O$-module, $\mathbf O$ et $\mathbf M$ les faisceaux
constants sur $X$ définis par $O$ et $M$. Il existe alors un foncteur spectral
cohomologique sur $\mathbf C^{\mathbf O}$, aboutissant au foncteur gradué
$(\mathrm{Ext}_{\mathbf O}^p(X;\mathbf M,A))$, et dont le terme initial est}
\LayerTarget{tohoku:eq:4.3.3}
\[
\tag{4.3.3}
\mathrm{II}_2^{p,q}(A)=\mathrm{Ext}_O^q(M,H^p(X,A)).
\]

\LayerTarget{tohoku:remark:edge-maps-and-five-term-sequence-4.3.1}
On en conclut comme d'habitude des ``edge-homomorphisms'' et une suite exacte
à 5 termes que le lecteur explicitera.
\LayerTarget{tohoku:construction:computing-constant-sheaf-spectral-functor}
Pour calculer le foncteur spectral, nous utiliserons le corollaire 1 à la
proposition 2.5.3, en utilisant le foncteur résolvant $\Gamma \mathrm C(A)$ de
$\Gamma$ (qu'on considère ici comme foncteur covariant de $\mathbf C^{\mathbf O}$
dans $\mathbf C^O$), $\mathrm C(A)$ étant la résolution canonique de $A$ (cf. 2.3), et le
foncteur résolvant $h_L$ de $h_M$, où $L$ est une résolution gauche de $M$ par
des $O$-modules libres. On trouve la première partie de la

\LayerTarget{tohoku:proposition:4.3.2}
\textsc{Proposition 4.3.2.} \textit{Sous les conditions du théorème
\textup{4.3.1}, choisissons une résolution gauche $L$ de $M$ par des $O$-modules
libres. Alors le foncteur spectral du théorème \textup{4.3.1} se calcule en prenant la
deuxième suite spectrale du bicomplexe $\mathrm{Hom}_O(L,\Gamma \mathrm C(A))$, où
$\mathrm C(A)$ est la résolution canonique de $A$ \textup{(}cf.\ \textup{2.3)}. Si on suppose les $L_i$
isomorphes à des $O^{n_i}$ \textup{(}$n_i$ entier $>0$\textup{)}, la première suite spectrale de
ce bicomplexe donne le foncteur spectral du théorème \textup{4.2.1}.}

\LayerTarget{tohoku:proof:proposition-4.3.2}
Pour la dernière assertion, il suffit de considérer la résolution $\mathbf L$
de $\mathbf M$ déduite de $L$, et de lui appliquer le corollaire 3 à la
proposition 4.2.4.
\LayerTarget{tohoku:reduction:kunneth-for-constant-sheaves}
La proposition précédente ramène le calcul de
$\mathrm{Ext}_{\mathbf O}(\mathbf M,A)$ et de ses suites spectrales à un
classique problème du type Künneth. Ainsi :

\TohokuContinuation{192}
